Prefix words modify the following instruction. One prefix word contains two 8-bit slots; the low byte is filled
and decoded first, and the high byte is filled and decoded second. Address-update and access-domain forms are
shown separately because address update is selected by effective-address syntax, while user access is selected
by memory operand annotation.

\subsection{Address-Update Operand Syntax}
Address update is no longer encoded as a prefix. Postincrement and predecrement are part of effective-address
syntax and are described in the Effective Addressing Modes chapter. The assembler spelling uses \texttt{++}
after the selected base or index term for postincrement and \texttt{{-}{-}} before the selected term for
predecrement.

@ADDRESS_UPDATE_TABLE@

\subsection{User Access Operand Syntax}
Access-domain prefixes annotate memory operands that should be interpreted in user or current access domains.
The prefix applies only to the immediately following instruction. These prefixes do not change the encoded EA
form; they select the access domain used after effective-address generation and address translation identify the
memory reference.

@ACCESS_DOMAIN_TABLE@

\subsection{Prefix Semantics}
A prefix word contains two independent 8-bit slots. Prefix bytes are interpreted from low byte to high byte and
modify only the immediately following instruction. If two prefixes in the same semantic group conflict, the later
decoded prefix overrides the earlier one. Unassigned prefix values have no architectural effect and remain
available for future definition.

@PREFIX_WORD_DIAGRAM@

@PREFIX_WORD_RULES@

@PREFIX_SLOT_TABLE@

\begin{description}[style=nextline,leftmargin=1.42in,labelwidth=1.32in,itemsep=3pt,topsep=3pt]
\item[NPX] Neutral prefix byte and canonical filler for an unused prefix slot.
\item[NOSPEC] Marks the following instruction as a speculation boundary or speculation-limited operation.
\item[SAT] Requests a saturating arithmetic variant for instruction forms that explicitly define saturation.
\item[NT] Provides a non-temporal memory hint without changing the architectural loaded or stored value.
\item[U2C] Selected by a user-domain source and current-domain destination memory annotation, for example
\texttt{MOV.Q U:[R0], R1} or \texttt{MOV.B U:[R0++], C:[R1++]}.
\item[C2U] Selected by a current-domain source and user-domain destination memory annotation, for example
\texttt{MOV.Q R0, U:[R1]} or \texttt{MOV.B C:[R0++], U:[R1++]}.
\item[U2U] Selected when all memory operands annotated by the instruction use the user access domain.
\end{description}

@REPCC_SECTION@

\subsection{Grouped Repeat Prefix REPG}
REPG repeats a byte-counted group of following instructions as a single architectural repeated operation. The
group length is explicit and must end on an instruction boundary.

Architectural execution remains scalar: each group iteration commits in program order, and faults are reported
at the instruction that faults within the current group iteration. The current encoding carries REPG as an
instruction form with an Rn counter and an immediate group byte count. This section defines the repeat semantics
because they are prefix-like at the execution-model level.

A REPG group is intended to describe short straight-line code. The group body must not contain a control
transfer into or out of the body, a nested repeat operation, an atomic read-modify-write operation, or a
state-changing system, TLB, cache, or fence operation unless a future instruction definition explicitly permits
that combination.

@REPG_RULE_TABLE@

@REPG_CLASS_TABLE@
